\icode{E-MCC-SI-L3-261023-001}
如图所示，质量为 $m$ 可看成质点的小球固定在一长为 $l$ 的轻杆上，轻杆另一端固定在底座的转轴上，底座的质量为 $M$，小球可绕转轴在竖直平面内自由转动。现在最低点给小球一个初速度，当小球通过最高点时，底座对地面的压力恰好为 $Mg$。在小球转动过程中，底座始终相对地面静止，不计转轴间摩擦阻力，重力加速度为 $g$，下列说法正确的是\paren[]

\begin{tikzpicture}
	\draw[dashed,cyan] (0,0) circle(1.5) (-1.5,0)--(1.5,0) (0,0)--(-135:1.9) (0,0)--(-90:1.9) (0,0)--(-45:1.9);
	\draw[fill=gray!50] (-45:1.9) arc(-45:-135:1.9)--(-135:3)--(-45:3)--(-45:1.9);
	\filldraw[pattern=north east lines] (-.05,0) rectangle(.05,1.5);
	\shade[ball color=gray] (0,1.5) circle(.1);
	\fill (0,0) circle(.05);
\end{tikzpicture}
\begin{choices}
	\item 外界对小球做的功为 $2mgl$
	\item 小球从最低点运动到最高点的过程中合外力的冲量大小为\((1 + \sqrt{5})mgl\)
	\item 小球运动到与转轴等高的位置时，地面对底座的静摩擦力大小为 $2mg$
	\item 小球在最低点时，底座对地面的压力大小为\((M + 5m)g\)
\end{choices}
\begin{answer}
    
\end{answer}



